\section*{第 \ref{ch:b3:planetary-atmospheres} 章——行星大气与宜居性}

\begin{solution}{exo:b3:planetary-atmospheres:1}
$H = RT/Mg$：金星 $8.314\times737/(0.04401\times8.87) = \qty{15.7}{km}$；
地球 $8.314\times288/(0.02897\times9.80) = \qty{8.4}{km}$；火星
$8.314\times214/(0.04334\times3.71) = \qty{11.1}{km}$；
$8.314\times94/(0.0280\times1.35) = \qty{20.7}{km}$。土卫六的大气，
尽管温度很低，却是最为膨胀的：其重力
比地球弱 7 倍。
\end{solution}

\begin{solution}{exo:b3:planetary-atmospheres:2}
$\lambda = GMm/kTr = \num{6.674e-11}\times\num{5.97e24}\times\num{1.674e-27}
/(\num{1.381e-23}\,T\times\num{6.878e6})$：在 \qty{800}{K} 时为 8.78，在
\qty{1600}{K} 时为 4.39。逃逸比例为 $9.78\,e^{-8.78} =
\num{1.5e-3}$ 和 $5.39\,e^{-4.39} = \num{6.7e-2}$：太阳活动周期
使氢的逃逸比例改变约 44 倍。
\end{solution}

\begin{solution}{exo:b3:planetary-atmospheres:3}
(a)~$d = S_\mathrm{eff}^{-1/2}$：\omterm{def:b3:planetary-atmospheres:runaway}{湿温室效应} \qty{0.993}{au}，\omterm{def:b3:planetary-atmospheres:runaway}{失控
温室效应} \qty{0.975}{au}，最大温室效应 \qty{1.705}{au}。(b)~金星位于 1.91，
接收的通量超过失控极限，甚至超过经验上的
近期金星极限（1.78）：它位于内边界以内。火星位于 0.43，
介于 1.014 与 0.3438 之间，位于保守宜居带之内：
其干燥源于它体积小以及大气丧失，
而非来自它的距离。
\end{solution}

\begin{solution}{exo:b3:planetary-atmospheres:4}
(a)~$\log_{10}p = 3.9895 - 443.028/93.66 = -0.741$, $p =
\qty{0.18}{bar}$。(b)~$0.05\times1.6 = \qty{0.080}{bar}$，相对湿度为
$0.080/0.18 = 44\,\%$。(c)~海面上方的空气远未饱和：
甲烷从海面蒸发，正如水从地球海洋中蒸发一样，
湿润空气须上升并冷却，然后才会降雨。
\end{solution}

\begin{solution}{exo:b3:planetary-atmospheres:5}
$\lambda_\mathrm D = 7.02\times2.014/1.008 = 14.0$. $u_\mathrm D/u_\mathrm H
= \sqrt{1.008/2.014} = 0.71$。每个原子的比值 $= 0.71\times(15.0\,
e^{-14.0})/(8.02\,e^{-7.02}) = 0.71\times\num{1.69e-3} = \num{1.2e-3}$。\omterm{def:b3:planetary-atmospheres:jeans-escape}{金斯逃逸}几乎不带走氘：
这是一种强筛选作用，
其 $\alpha \approx 0.001$。
\end{solution}

\begin{solution}{exo:b3:planetary-atmospheres:6}
(a)~$0.016/\num{1.5575e-4} = 103$。(b)~当 $\alpha = 0$ 时，富集倍数为
$M_0/M$：$M_0 = 103\times\num{6e15} = \qty{6.2e17}{kg}$，或
$\num{6.2e17}/(1000\times\num{4.60e14}) = \qty{1.3}{m}$ 厚的水层覆盖
金星。(c)~$M_0/M = 103^{1/0.2} = \num{1.2e10}$，过去的存量
为 \qty{7e25}{kg}，是地球质量的十二倍：不可能。
富集需要强分馏逃逸（\omterm{def:b3:planetary-atmospheres:jeans-escape}{金斯逃逸}或\omterm{def:b3:planetary-atmospheres:hydrodynamic}{非热逃逸}），
至少在最后一部分水流失期间发挥作用，或者处于
有水补给的稳态。
\end{solution}

\begin{solution}{exo:b3:planetary-atmospheres:7}
(a)~在四十亿年（$\num{1.26e17}$~s）内火星损失
$\num{3.8e43}$ 个原子，$\num{1.9e43}$ 个水分子，$\qty{5.7e17}{kg}$；
在 $4\pi(\num{3.3895e6})^2 = \qty{1.44e14}{m^2}$ 的面积上，相当于一层
厚 \qty{3.9}{m} 的水层。(b)~多出二十倍，达 \qty{78}{m}。(c)~即使是高倾角时的速率，
也只达到 100 至
\qty{1500}{m} 范围的下限；其余的水进入了地壳与冰层，
或在太阳年轻时流失得更快。
\end{solution}

\begin{solution}{exo:b3:planetary-atmospheres:8}
(a)~$m = pA/g = 636\times\num{1.44e14}/3.71 = \qty{2.5e16}{kg}$.
(b)~$0.1\times\num{1.26e17} = \qty{1.3e16}{kg}$，相当于现今大气的一半；
全部耗尽所需时间为 $\num{2.5e16}/0.1 = \qty{2.5e17}{s}$，
\num{7.8}~十亿年。(c)~年轻的太阳更为活跃，风暴更为频繁，
流失在风暴期间增强；其他途径
（氢逃逸、溅射、光化学逃逸）未计入。
\end{solution}

\begin{solution}{exo:b3:planetary-atmospheres:9}
(a)~由 $S(1 - 0.193)/4 = 291$ 得 $S = \qty{1442}{W.m^{-2}}$，为
地球的 1.06 倍。(b)~今日金星将吸收 $0.807\times2601/4 =
\qty{525}{W.m^{-2}}$；年轻的金星，$0.7\times2601 = \qty{1821}{W.m^{-2}}$
接收，吸收 \qty{367}{W.m^{-2}}：两者都超过极限，因此无云的年轻金星
从一开始就已失控。(c)~云会提高反照率；
三维含云模式已找到
通往宜居年轻金星的途径。
\end{solution}

\begin{solution}{exo:b3:planetary-atmospheres:10}
(a)~$H = kT/mg = \num{1.381e-23}\times1000/(16\times\num{1.661e-27}\times9.21)
= \qty{56}{km}$。(b)~$n_\mathrm{exo} = 1/(\sigma_c H) = 1/(\num{3e-19}\times
\num{5.64e4}) = \qty{5.9e13}{m^{-3}}$; $z = 200 + 56.4\ln(\num{1e16}/
\num{5.9e13}) = 200 + 56.4\times5.13 = \qty{490}{km}$。(c)~高处的重力较弱，
使标高随高度增大，因此密度下降
比假定的更慢：\omterm{def:b3:planetary-atmospheres:exobase}{外逸层底}上移。
\end{solution}

\begin{solution}{exo:b3:planetary-atmospheres:11}
(a)~$\Phi = \num{3.0e13}\times\num{3e-3} = \qty{9e10}{cm^{-2}.s^{-1}} =
\qty{9e14}{m^{-2}.s^{-1}}$；在 $4\pi(\num{6.371e6})^2 =
\qty{5.10e14}{m^2}$，每秒共 $\num{4.6e29}$ 个 $\ce{H2}$ 分子。
(b)~海洋含有 $\num{1.4e24}/18.015\times\num{6.022e23} =
\num{4.7e46}$ 个水分子，每个对应一个 $\ce{H2}$：其持续时间为
$\qty{1.0e17}{s} = \num{3.2}\text{ 十亿年}$，与
地球年龄相当，这就是为什么混合比 \num{3e-3} 标志着\omterm{def:b3:planetary-atmospheres:runaway}{湿温室
效应}极限。
\end{solution}

\begin{solution}{exo:b3:planetary-atmospheres:12}
(a)~$p_1 = g/D\kappa = 9.81/(1.66\times0.005) = \qty{1180}{Pa}$.
(b)~$1/T = 1/273.16 - \ln(1182/611.657)\times461.5/\num{2.5e6}$, $T =
\qty{282.5}{K}$，且 $\sigma T^4 = \qty{361}{W.m^{-2}}$。(c)~真实极限
更低，因此辐射层更冷、气压更低：真实
水汽比这种灰气体更不透明，其有效 $\kappa$
更大。
\end{solution}

\begin{solution}{pb:b3:planetary-atmospheres:1}
\textbf{1.} $g = GM/R^2 = \num{6.674e-11}\times\num{5.2e24}/(\num{6.2e6})^2 =
\qty{9.03}{m.s^{-2}}$；$\bar\rho = M/\tfrac43\pi R^3 = \qty{5210}{kg.m^{-3}}$。

\textbf{2.} $v_\mathrm{esc} = \sqrt{2GM/r} = \sqrt{2\times\num{3.47e14}/\num{6.7e6}}
= \qty{10.18}{km.s^{-1}}$。

\textbf{3.} $u = \sqrt{2kT/m}$：氢 \qty{4.81}{km.s^{-1}}，氦
\qty{2.41}{km.s^{-1}}，氧 \qty{1.21}{km.s^{-1}}。

\textbf{4.} $\lambda = (v_\mathrm{esc}/u)^2$：4.49，17.8，71.2。

\textbf{5.} 只有氢：其逃逸比例为 0.06，氦为
\num{3.5e-7}，氧为 \num{9e-30}。三者都不接近流体动力学区，
$\lambda \approx 1.5$ 至 3。

\textbf{6.} $\lambda = 4.49\times1400/3000 = 2.09$：处于过渡区；
氢本会以流体动力学方式逸出，
或许还带走较重的气体。

\textbf{7.} $(1 + 4.49)\,e^{-4.49} = 0.0618$。

\textbf{8.} $n u/2\sqrt\pi = \num{2e11}\times4806/3.545 = \qty{2.71e14}{m^{-2}.s^{-1}}$；
$\Phi_\mathrm J = \num{2.71e14}\times0.0618 = \qty{1.68e13}{m^{-2}.s^{-1}}$。

\textbf{9.} $4\pi(\num{6.7e6})^2 = \qty{5.64e14}{m^2}$；每秒
$\num{9.5e27}$ 个原子。

\textbf{10.} $\num{1.4e24}/18.015 = \qty{7.8e22}{mol}$，$\num{4.68e46}$ 个分子，
$\num{9.36e46}$ 个氢原子。

\textbf{11.} $\num{9.36e46}/\num{9.5e27} = \qty{9.9e18}{s}$，约 3100 亿
年。

\textbf{12.} $\lambda_\mathrm D = 4.49\times2.014/1.008 = 8.96$；逃逸比例
$9.96\,e^{-8.96} = \num{1.28e-3}$；相对效率 $0.707\times
\num{1.28e-3}/0.0618 = 0.015$。

\textbf{13.} $H_a = kT/m_a g = \num{1.381e-23}\times200/(44.01\times\num{1.661e-27}
\times9.03) = \qty{4.18}{km}$。

\textbf{14.} $b = \qty{1.2e21}{m^{-1}.s^{-1}}$，于是
\[
\Phi_\mathrm{DL} = \frac{b f(m_a - m_i)}{m_a H_a}
= \frac{\num{1.2e21}\times\num{1e-3}\times(42.0/44.0)}{4185}
= \qty{2.74e14}{m^{-2}.s^{-1}}.
\]

\textbf{15.} $\num{3.0e13}\times\num{1e-3}\times9.03/9.81 =
\qty{2.76e10}{cm^{-2}.s^{-1}} = \qty{2.76e14}{m^{-2}.s^{-1}}$：两者在
1\,\% 之内相同。

\textbf{16.} $4\pi(\num{6.2e6})^2 = \qty{4.83e14}{m^2}$；$\num{1.32e29}$ 个
$\ce{H2}$ 分子，每秒 $\num{2.64e29}$ 个原子。

\textbf{17.} 扩散限制所供应的量，是\omterm{def:b3:planetary-atmospheres:exobase}{外逸层底}
在当前密度下放出量的 28 倍。氢在热层中堆积，直至
金斯通量与供应量相等：长期来看，逃逸受扩散限制，
与地球上的情形相同。

\textbf{18.} $\num{2e11}\times28 = \qty{5.6e12}{m^{-3}}$。

\textbf{19.} $\num{4.68e46}/\num{1.32e29} = \qty{3.54e17}{s}$，112 亿年。

\textbf{20.} $S = 0.30/0.60^2 = 0.833$。

\textbf{21.} $d = \sqrt{L/S_\mathrm{eff}}$：内缘 $\sqrt{0.30/0.939} =
\qty{0.565}{au}$，外缘 $\sqrt{0.30/0.290} = \qty{1.02}{au}$。
Kepos 位于 \qty{0.60}{au}，处于内侧，靠近内缘。

\textbf{22.} $0.939/0.833 = 1.13$：增亮 13\,\%。

\textbf{23.} $0.833\times1361\times0.7/4 = \qty{198}{W.m^{-2}}$，$T_\mathrm{eq} =
(198/\sigma)^{1/4} = \qty{243}{K}$；远低于
\qty{291}{W.m^{-2}} 的极限（按类太阳恒星计算，
但余量仍很大）。

\textbf{24.} 快三倍：37 亿年。

\textbf{25.} \textbf{110~亿年}。
\end{solution}
